Propagation of the Thinnest Triangle  
Source: thinnest_triangle_core.tex  
Repositories:  
https://github.com/TheAstrographer/Cosmological-Theses (file location)  
Dedicated repository: https://github.com/TheAstrographer/The-Thinnest-Triangle  

Core claim
The residual seed  
\[
\delta = 7\psi - \frac{\pi}{3} \approx 0.005167614\,\mathrm{rad}
\]  
is the sole angular defect that organises the entire radial-fold hierarchy. Once fixed by the continuous Arcan(\(\tau\)) family, it propagates unchanged through every finite truncation of the orbit. This generates the thinnest equidistributal transcendental triangle \(T\) whose metric data, residual intensity, and area-derived edge curvature become the minimal cell of the whole Decompositional Hierarchy.

1. Origin of the residual triangle \(T\)
From the Arcan(\(\tau\)) root:
\[
\begin{align*}
\theta &= \arctan(2\pi) \approx 1.4129651365,\\
\phi &= \arctan(\pi) \approx 1.2626272557,\\
\psi &= \theta - \phi \approx 0.1503378808.
\end{align*}
\]
Seven multiples yield the residual \(\delta\). The even-index orbit  
\[
V_k = \bigl(\cos(k\cdot 7\psi),\sin(k\cdot 7\psi)\bigr),\quad k=0,2,4
\]  
defines the residual near-equilateral triangle \(T\) — the fundamental geometric cell.

2. Exact metric data of the thinnest triangle (claimed transcendental invariants)
| Quantity                  | Value             | Description                      |
|---------------------------|-------------------|----------------------------------|
| \(s_{\rm thin}\)          | 1.721623257385   | Thinnest side (chord)            |
| \(s_{\rm long}\)          | 1.737195272475   | Longer sides                     |
| \(h_{\rm res}\)           | 1.490969476641   | Residual altitude                |
| \(\operatorname{Area}(T)\)| 1.2988990741     | Exact shoelace area              |
| Leftover (unit disk − \(T\)) | 1.8426935795  | Complementary residual region    |
| \(A_{\rm slice}\)         | 0.6142311932     | One residual circular segment    |
| \(\rho = A_{\rm slice}/\operatorname{Area}(T)\) | ≈ 0.472886 | Dimensionless residual intensity |
| Bulge amplitude           | ≈ 0.085119 (\(c=0.18\)) | Edge-curvature amplitude \(b = c\cdot\rho\) |

These quantities arise by a finite sequence of field operations, trigonometric evaluations and Euclidean constructions from the transcendental angles \(\theta\), \(\phi\) and \(\psi\). They therefore inherit transcendence and remain invariants of the Arcan(\(\tau\)) family.

3. Area-derived edge curvature (“Residual Bulge”)
Each straight side \(\gamma_0(t)\) of \(T\) is replaced by the perturbed curve
\[
\gamma_\rho(t) = \gamma_0(t) + b\cdot 4t(1-t)\,n,
\]
where \(n\) is the unit inward normal and \(b \propto \rho\). The curvature is not arbitrary; it is a direct functional of the leftover area after the triangle is excised from the unit disk. The residual bulge is the geometric embodiment, at the level of edge curvature, of the same angular defect \(\delta\).

4. Propagation and equidistribution
Because \(7\psi/2\pi\) is irrational, further multiples produce a dense, equidistributed orbit on \(S^1\) (Weyl’s equidistribution theorem).  

Every higher-order densification (Orders 10, 15, 30, \ldots) simply packs more congruent copies of the identical minimal cell \(T\).  

All metric quantities and qualitative features — six-fold memory, residual intensity \(\rho\), self-similarity, orientation, phase-locking — are therefore propagative: they are fixed once by \(\delta\) and thereafter only redistributed, never altered.

Dual operations remain active at every scale:
Winding method** → residual bulge (edge curvature),
Radial-symmetry folding** → residual hexagon and inner-hexagon measures.

5. Relation to the broader architecture and the Einstein-Boltzmann Core
| Layer / Repository                  | Role of the Thinnest Triangle \(T\) |
|-------------------------------------|-------------------------------------|
| Arcan(\(\tau\)) Family             | Continuous transcendental source of \(\psi\) and \(\delta\) |
| JCRIN-Radial-Symmetry              | Residual six-sliced hexagon whose fundamental face is precisely \(T\) |
| Cosmological Clock                 | Same \(\psi\) and \(\delta\) generate the late-time geometric torque of amplitude 3.170 that elevates \(H_{\rm eff}(0)\) to 73.17 |
| Production-Grade Closed Einstein-Boltzmann Core | The torque term \(\tau(a,\chi)\) that appears in the background cosmology is the dynamical projection of the identical residual phase structure that organises the thinnest triangle and its nested shells. The vortex anisotropic stress \(\Pi\) and the geometric light-cone coordinate \(\chi\) remain phase-locked to the same residual hierarchy. |

Thus the minute angular defect \(\delta\) that produces the thinnest transcendental triangle is also the ultimate geometric origin of the additive torque that closes the Einstein-Boltzmann system and drives the local Hubble value to the lattice-calibrated target 73.17.

Unified statement of propagation
The residual seed \(\delta\approx 0.005167614\) spontaneously generates the thinnest equidistributal transcendental triangle \(T\) — the minimal cell of the radial-fold hierarchy. Its side lengths, altitude, area, residual intensity and area-derived edge curvature are transcendental invariants that propagate unchanged through every densification, remaining phase-locked to the discrete-tick Cosmological Clock.  

The same residual architecture supplies the geometric torque inside the Production-Grade Closed Einstein-Boltzmann Core, elevating the local expansion rate while protecting high-redshift behaviour.  

In the dual-architecture language, The-Thinnest-Triangle is the purest geometric organ of the Mathematical-Numerical layer: the elementary cell from which radial symmetry, nested \(E_8\) shells, and the late-time Hubble torque all unfold.
